\section{Conclusion}
\label{sec:conclusion}

\TW{} starts from the premise that an AEAD tag should be the compact digest
applications often treat it as: a hash output for the full encryption context,
with committing security following from that hash story rather than added as an
external layer. The construction adapts the KangarooTwelve tree-hashing pattern
into a two-level tree of keyed duplexes over a single \Keccakp{} permutation.
The chunk decomposition is canonical, while the number of leaf duplexes
evaluated in parallel is an implementation choice rather than a mode parameter.
This gives scalar, SIMD, and multicore implementations the same ciphertexts and
tags, with constant working memory and scalable parallelism.

The main security result is tag hash security for the wrapper that returns the
final authentication tag on input $(K, U, A, M)$. In the public permutation
model, via an intermediate random oracle and the flat sponge lift, we prove
collision resistance, standard and everywhere second-preimage resistance,
preimage resistance on fixed input slices, and everywhere preimage resistance.
Those bounds yield the AEAD-facing committing statements, including full-input
\textsf{CMTDs-4}. The same root structure gives \textsf{INT-RUP}
integrity under release of unverified plaintext. As baseline AEAD properties, we
prove concrete multi-user indistinguishability and derive all-in-one AE as its
single-user corollary. At the implemented parameters, all stated notions meet a
$128$-bit target.

The implementation results show the cost and payoff of putting that tag hash
inside a tree. Short inputs pay for the larger permutation and the root/leaf
structure, but long inputs expose independent leaves while preserving one
canonical wire format. Vectorized \texttt{amd64} and \texttt{arm64}
implementations reach bulk throughputs of $48.0$\,Gbit/s on the AVX-512 path,
$39.4$\,Gbit/s on the \texttt{arm64} path, and $16.6$\,Gbit/s on the AVX2
path, with the AVX-512 and \texttt{arm64} paths at $51$--$76\%$ of
hardware-accelerated AES-128-GCM throughput.

The remaining questions are those left open in \Cref{sec:limitations}: nonce
misuse would require a different construction, privacy under release of
unverified plaintext is not provided by the \textsf{INT-RUP} theorem, and sharper
bounds would need either a direct keyed-duplex treatment of the BDPV11-style
phasing with adversarial initialization or a beyond-birthday argument. More
broadly, \TW{} supports the thesis that a single \Keccak{} permutation with
substantial cryptanalytic history can provide authenticated encryption whose tag
has hash security strong enough to drive commitment, while retaining scalable
hashing-style performance.
